\documentclass[11pt]{article}

\usepackage[T1]{fontenc}
\usepackage{lmodern}
\usepackage[margin=1in]{geometry}
\usepackage{amsmath,amssymb,amsthm}
\usepackage{booktabs}
\usepackage{microtype}
\usepackage[hidelinks]{hyperref}

\newtheorem{theorem}{Theorem}[section]
\newtheorem{proposition}[theorem]{Proposition}
\newtheorem{lemma}[theorem]{Lemma}

\newcommand{\A}{\mathbb A}
\newcommand{\N}{\mathbb N}
\newcommand{\Lang}{\mathcal L}
\newcommand{\one}{\mathbf 1}

\hypersetup{
  pdftitle={A Ternary Word Attaining the Abelian Maximal Pattern Complexity Bound at Pattern Size Three},
  pdfauthor={Carptopus},
  pdfsubject={Abelian maximal pattern complexity of a ternary substitution word},
  pdfkeywords={Abelian maximal pattern complexity, primitive substitution, aperiodic by projection, computer-assisted proof}
}

\title{A Ternary Word Attaining the Abelian Maximal Pattern\\
Complexity Bound at Pattern Size Three}
\author{Carptopus\\\small\texttt{carptopus@163.com}}
\date{25 August 2026\\\small v0.1-beta; not peer reviewed}

\begin{document}
\maketitle

\begin{abstract}
Kamae, Widmer and Zamboni proved that every recurrent word over an alphabet of
size $r$ which is aperiodic by projection has Abelian maximal pattern
complexity at least $(r-1)k+1$ at pattern size $k$. Their introduction left
open whether this bound is attainable when $r\ge3$ and $k\ge3$. We settle the
smallest unresolved parameter pair under the literature search described
below. Let $\alpha$ be the fixed point beginning in $0$ of the constant-length
substitution
\[
0\mapsto001,\qquad 1\mapsto020,\qquad 2\mapsto000.
\]
We prove that $\alpha$ is recurrent and aperiodic by projection and that
\[
p_{\alpha}^{*\mathrm{ab}}(3)=7.
\]
The upper bound is computer-assisted but exact. Sparse three-point languages
obey a residue recursion modulo $3$. A $3\times3$ patch of neighboring
languages makes this recursion finite-state; the closure contains 938 states
and covers every pair of nonnegative offsets through their base-$3$ digits.
Every central state has at most seven Parikh vectors, while the pattern
$\{0,2,9\}$ realizes seven. The verification program is supplied with a
canonical transition-table hash and fail-closed checks.
\end{abstract}

\noindent\textbf{Keywords.} Abelian maximal pattern complexity; infinite word;
primitive substitution; aperiodic by projection; automatic sequence;
computer-assisted proof.

\medskip
\noindent\textbf{2020 Mathematics Subject Classification.} 68R15, 37B10.

\section{Introduction}

Let $\A$ be a finite alphabet and let
$\alpha=\alpha_0\alpha_1\alpha_2\cdots\in\A^{\N}$. For a finite set
$S=\{s_1<\cdots<s_k\}\subset\N$, an $S$-factor is
\[
\alpha[n+S]=\alpha_{n+s_1}\alpha_{n+s_2}\cdots\alpha_{n+s_k}.
\]
Two finite words are Abelian equivalent when they have the same Parikh vector.
Write $p_\alpha^{\mathrm{ab}}(S)$ for the number of Abelian classes among all
$S$-factors and
\[
p_\alpha^{*\mathrm{ab}}(k)
=\sup_{|S|=k}p_\alpha^{\mathrm{ab}}(S).
\]

We call $\alpha$ recurrent if each of its finite factors occurs infinitely
often. For a nonempty proper subset $B\subsetneq\A$, let
$\one_B(\alpha)$ be the binary word whose $n$-th term records whether
$\alpha_n\in B$. The word $\alpha$ is \emph{aperiodic by projection} if
$\one_B(\alpha)$ is not eventually periodic for every such $B$.

Kamae, Widmer and Zamboni~\cite[Theorem~4]{KWZ} proved that if $\alpha$ is
recurrent and aperiodic by projection over an $r$-letter alphabet, then
\begin{equation}
p_\alpha^{*\mathrm{ab}}(k)\ge(r-1)k+1.
\label{eq:kwz}
\end{equation}
For binary words equality is automatic, and for arbitrary $r$ they gave
equality examples at $k=2$. Their introduction records the pointwise
sharpness problem for $r\ge3,k\ge3$. Their Remark~4.2 asks the stronger and
still separate question whether one word can attain equality for every $k$.
The survey of Fici and Puzynina~\cite{FP} lists only the previously known
equality cases. The present note settles $r=3,k=3$, the smallest unresolved
parameter pair under our current search; it does not answer the stronger
all-$k$ question.

Our construction is a primitive constant-length substitution rather than the
2-adic Toeplitz family used in~\cite{KWZ} for $k=2$. The main difficulty is not
finding seven classes for one pattern, but proving that no three-point pattern
has eight. A bounded scan over offsets cannot establish that assertion. We
instead give an exact digit automaton whose finite closure covers all offset
pairs.

\section{Main result}

Define the length-three substitution $\sigma$ on $\A=\{0,1,2\}$ by
\begin{equation}
\sigma(0)=001,\qquad \sigma(1)=020,\qquad \sigma(2)=000.
\label{eq:substitution}
\end{equation}
Since $\sigma(0)$ begins with $0$, the limit
\[
\alpha=\lim_{m\to\infty}\sigma^m(0)
\]
is a well-defined one-sided fixed point.

\begin{theorem}\label{thm:main}
The word $\alpha$ is recurrent and aperiodic by projection, and
\[
\boxed{p_\alpha^{*\mathrm{ab}}(3)=7.}
\]
Thus the lower bound~\eqref{eq:kwz} is sharp for $r=3,k=3$.
\end{theorem}

The proof occupies Sections~\ref{sec:dynamics}--\ref{sec:certificate}.
Section~\ref{sec:dynamics} verifies the dynamical hypotheses,
Section~\ref{sec:recursion} gives the exact sparse-language recursion, and
Section~\ref{sec:certificate} converts that recursion into a finite
certificate covering every pattern.

\section{Recurrence and aperiodicity by projection}\label{sec:dynamics}

\begin{proposition}\label{prop:recurrence}
The substitution $\sigma$ is primitive. Consequently its fixed point
$\alpha$ is uniformly recurrent.
\end{proposition}

\begin{proof}
The word $\sigma(0)$ contains $0,1$, $\sigma(1)$ contains $0,2$, and
$\sigma(2)$ contains $0$. Iterating at most three further times therefore
makes every letter occur in the image of every letter. Hence the incidence
matrix of $\sigma$ is primitive. The standard primitive-substitution theorem
\cite{Queffelec} implies uniform recurrence of the fixed point.
\end{proof}

For $i\in\{0,1,2\}$, put
\[
\chi_i(n)=\one_{\{\alpha_n=i\}}.
\]
Reading~\eqref{eq:substitution} coordinatewise gives
\begin{align}
\alpha_{3n}&=0,\nonumber\\
\chi_2(3n+1)&=\chi_1(n),\nonumber\\
\chi_1(3n+2)&=\chi_0(n),\nonumber\\
\chi_0(3n+2)&=1-\chi_0(n).
\label{eq:digit-identities}
\end{align}
Moreover, $\chi_2$ is supported on $1\pmod3$, while $\chi_1$ is supported on
$2\pmod3$.

\begin{proposition}\label{prop:projection}
The word $\alpha$ is aperiodic by projection.
\end{proposition}

\begin{proof}
Primitivity implies that each letter occurs infinitely often. Suppose first
that $\chi_0$ is eventually periodic with positive period $p$. If $3\nmid p$,
the identity $\chi_0(3n)=1$, together with $\gcd(3,p)=1$, forces $\chi_0$
eventually to be identically one. This contradicts the infinite occurrence of
$1$ and $2$. If $p=3q$, then the last identity in
\eqref{eq:digit-identities} shows that $q$ is also an eventual period of
$\chi_0$. Removing all powers of $3$ from $p$ reduces to the preceding
contradiction.

If $\chi_1$ is eventually periodic, its support in one residue class modulo
$3$, together with infinitely many ones, forces every eventual period to be
divisible by $3$. The third identity in~\eqref{eq:digit-identities} then
transfers an eventual period divided by $3$ to $\chi_0$, which is impossible.
The same argument using the second identity transfers eventual periodicity of
$\chi_2$ first to $\chi_1$ and then to $\chi_0$.

Every nonempty proper subset of a three-letter alphabet is either a singleton
or the complement of a singleton. Its indicator is therefore one of the
$\chi_i$, or its complement. No nontrivial projection is eventually periodic.
\end{proof}

\section{Exact recursion for sparse languages}\label{sec:recursion}

For an ordered integer tuple $C=(c_1,\ldots,c_m)$, define
\[
\Lang(C)=\{(\alpha_{n+c_1},\ldots,\alpha_{n+c_m}):
\text{the tuple occurs for arbitrarily large }n
\text{ with }n+c_i\ge0\text{ for all }i\}.
\]
Adding the same integer to every coordinate does not change $\Lang(C)$.
Uniform recurrence also ensures that deleting a finite initial segment does
not change this language. Let
\[
f_d(x)=\sigma(x)_d,\qquad d\in\{0,1,2\}.
\]

\begin{lemma}[residue recursion]\label{lem:recursion}
Normalize $C$ by subtracting $\min_i c_i$. For each
$t\in\{0,1,2\}$, put
\[
q_i=\left\lfloor\frac{t+c_i}{3}\right\rfloor,
\qquad d_i\equiv t+c_i\pmod3,
\]
and normalize $Q_t=(q_i)_i$ by a common translation. Then
\begin{equation}
\Lang(C)=\bigcup_{t=0}^2
\{(f_{d_1}(x_1),\ldots,f_{d_m}(x_m)):
(x_1,\ldots,x_m)\in\Lang(Q_t)\}.
\label{eq:recursion}
\end{equation}
\end{lemma}

\begin{proof}
Write the translation origin as $n=3u+t$. Each sampled coordinate then lies
in position $d_i$ of the substituted letter $\alpha_{u+q_i}$. Conversely,
every occurrence of a child tuple at an arbitrarily large $u$ produces an
occurrence in the corresponding residue class. Common normalization only
shifts $u$, and recurrence removes the finite boundary. No coordinate
reordering is performed.
\end{proof}

The recursion terminates at span one. The exact two-letter factor set is
\begin{equation}
\mathcal F_\alpha(2)=\{00,01,02,10,20\}.
\label{eq:pairs}
\end{equation}
Indeed, internal factors of the three substitution images give
$00,01,02,20$; every image starts in $0$, while its final letter is $0$ or
$1$, so crossing a substitution boundary adds only $00$ and $10$.

\section{The finite patch certificate}\label{sec:certificate}

For integers $a,b$, define the nine-cell patch
\begin{equation}
K(a,b)=\big(\Lang(0,a+i,b+j)\big)_{i,j\in\{-1,0,1\}}.
\label{eq:patch}
\end{equation}

\begin{lemma}[digit transition]\label{lem:transition}
For every $x,y\in\{0,1,2\}$, there is an explicit deterministic map
$\Phi_{x,y}$ such that
\begin{equation}
K(3a+x,3b+y)=\Phi_{x,y}(K(a,b)).
\label{eq:transition}
\end{equation}
Write $K_{u,v}$ for the $(u,v)$-cell of a patch. For each target cell
$(i,j)\in\{-1,0,1\}^2$ and $t\in\{0,1,2\}$, set
\[
u_t=\left\lfloor\frac{t+x+i}{3}\right\rfloor,
\qquad
v_t=\left\lfloor\frac{t+y+j}{3}\right\rfloor,
\]
and
\[
d_t=(t,\ (t+x+i)\bmod3,\ (t+y+j)\bmod3).
\]
Then
\begin{align}
[\Phi_{x,y}(K)]_{i,j}
=\bigcup_{t=0}^2\{&(
f_{d_{t,0}}(z_0),f_{d_{t,1}}(z_1),f_{d_{t,2}}(z_2)):\nonumber\\[-2mm]
&(z_0,z_1,z_2)\in K_{u_t,v_t}\}.
\label{eq:explicit-transition}
\end{align}
\end{lemma}

\begin{proof}
Apply~\eqref{eq:recursion} to each of the nine target cells. After division
by $3$, each quotient offset differs from $a$ or $b$ by an element of
$\{-1,0,1\}$. Thus every required child language is already a cell of
$K(a,b)$; the output digits determine the coordinate maps $f_d$.
\end{proof}

Starting from $K(0,0)$, close under the nine maps $\Phi_{x,y}$. The fail-closed
verification obtains the following certificate.

\begin{theorem}[finite closure certificate]\label{thm:certificate}
The transition closure contains 938 patch states and 8442 directed labelled
transitions. For the central language of a state, the distribution of the
number of Parikh vectors is
\begin{center}
\begin{tabular}{@{}rr@{}}
\toprule
number of vectors & number of states\\
\midrule
3 & 3\\
4 & 246\\
5 & 223\\
6 & 285\\
7 & 181\\
\bottomrule
\end{tabular}
\end{center}
The maximum is seven. The canonical JSON serialization of the states and all
labelled transitions has SHA-256 digest
\begin{center}
\small\texttt{BBF0CBF7D6E5E23CCEECC71A59644AD4F5DA0D48A1A2A0A219BB292A127CDAAA}.
\end{center}
\end{theorem}

\begin{proof}[Why this covers every pattern]
Let $0\le a,b$. Read the paired base-$3$ digits of $a,b$ from most significant
to least significant, padding with leading zeros. Repeatedly apply the
corresponding transition $\Phi_{x,y}$ to $K(0,0)$. Equation
\eqref{eq:transition} shows inductively that the final state is $K(a,b)$.
Therefore this finite closure covers every nonnegative offset pair; it is not
a bounded search over offsets. In particular, every three-point pattern
$\{s_1<s_2<s_3\}$, after translation to
$\{0,s_2-s_1,s_3-s_1\}$, has at most seven Abelian classes.

The supplied program cross-checks the Parikh images of the recursive languages
against a length $3^{10}$ fixed-point prefix for a small offset box, checks
digit transitions against direct recursive patches for $0\le a,b<100$, and
aborts unless all state counts, the distribution and the canonical digest
match the frozen values.
\end{proof}

\begin{proposition}[sharp pattern]\label{prop:witness}
For $S=\{0,2,9\}$, the exact Parikh set is
\[
\{(0,2,1),(1,0,2),(1,1,1),(1,2,0),
(2,0,1),(2,1,0),(3,0,0)\}.
\]
Hence $p_\alpha^{\mathrm{ab}}(S)=7$.
\end{proposition}

Combining Theorem~\ref{thm:certificate} and Proposition~\ref{prop:witness}
proves the complexity equality in Theorem~\ref{thm:main}. Propositions
\ref{prop:recurrence} and~\ref{prop:projection} verify the hypotheses of
\eqref{eq:kwz}, providing an independent lower-bound route as well.

\section{Scope and prior-work boundary}

The result establishes sharpness only at $r=3,k=3$. It does not prove that the
KWZ bound is sharp for all $r\ge3,k\ge3$, and it does not produce a single word
attaining the bound for every $k$.

The candidate new contribution is the explicit ternary substitution and the
all-pattern proof that its three-point Abelian maximal pattern complexity is
seven. Primitive-substitution recurrence and the KWZ lower bound are standard
inputs. The finite-state method is used here to certify the global upper
bound; the program is part of the proof package rather than evidence from a
finite sample.

The literature search covered the original paper~\cite{KWZ}, the
survey~\cite{FP}, visible citation metadata, and targeted searches for the
substitution and the exact parameter pair. No direct prior construction was
found. This is not an absolute-priority claim: unindexed work, different
notation and concurrent research remain possible until broader expert review.

\section{Reproducibility and responsibility}

The archived verification file is
\texttt{verification/verify\_ternary\_ampc\_k3.py}, with SHA-256
\begin{center}
\small\texttt{7F827D487681F3F6139FBB239E89EAD0A5303C6CE853866B6F6E6E082CE78CC8}.
\end{center}
It uses Python 3.11 or later and only the standard library. From the entry
directory, run
\begin{center}
\texttt{python -X utf8 verification/verify\_ternary\_ampc\_k3.py}.
\end{center}
The program regenerates the finite closure, checks the frozen counts and
closure hash, and prints the sharpness witness. No random seed or external
solver is used. The verification README records the expected output, file
manifest and evidence boundary.

Carptopus is the responsible author and takes responsibility for the
mathematical claims and the released materials. OpenAI Codex was used
extensively for AI-assisted research, proof development, verification and
writing. The result remains subject to public mathematical review.

\begin{thebibliography}{9}

\bibitem{KWZ}
T.~Kamae, S.~Widmer and L.~Q.~Zamboni,
\newblock Abelian maximal pattern complexity of words,
\newblock \emph{Ergodic Theory Dynam. Systems} 35 (2015), 142--151.
\newblock \href{https://doi.org/10.1017/etds.2013.51}{doi:10.1017/etds.2013.51}.

\bibitem{FP}
G.~Fici and S.~Puzynina,
\newblock Abelian combinatorics on words: a survey,
\newblock \emph{Comput. Sci. Rev.} 47 (2023), 100532.
\newblock \href{https://doi.org/10.1016/j.cosrev.2022.100532}{doi:10.1016/j.cosrev.2022.100532}.

\bibitem{Queffelec}
M.~Queffelec,
\newblock \emph{Substitution Dynamical Systems---Spectral Analysis},
\newblock second ed., Lecture Notes in Mathematics 1294, Springer, 2010.

\end{thebibliography}

\end{document}
